%%%%%%%%%%%%%%%%%%%%%%%%%%%%%%%
%%%Autor: Briggette Román
%%%%%%%%%%%%%%%%%%%%%%%%%%%%%%%


\documentclass[12pt,a4paper]{book}
%%%%%%%%%%%%%%%%%
%Manejo de paquetes
%%%%%%%%%%%%%%%%%
\usepackage[utf8]{inputenc} 
\usepackage[spanish]{babel} 
\usepackage{hyperref}
\usepackage{algorithm}
\usepackage{algorithmic}
\floatname{algorithm}{Algoritmo}
\renewcommand{\listalgorithmname}{Lista de algoritmos}
\renewcommand{\algorithmicrequire}{\textbf{Entrada:}}
\renewcommand{\algorithmicensure}{\textbf{Salida:}}
\renewcommand{\algorithmicend}{\textbf{fin}}
\renewcommand{\algorithmicif}{\textbf{si}}
\renewcommand{\algorithmicthen}{\textbf{entonces}}
\renewcommand{\algorithmicelse}{\textbf{si no}}
\renewcommand{\algorithmicelsif}{\algorithmicelse,\ \algorithmicif}
\renewcommand{\algorithmicendif}{\algorithmicend\ \algorithmicif}
\renewcommand{\algorithmicfor}{\textbf{para}}
\renewcommand{\algorithmicforall}{\textbf{para todo}}
\renewcommand{\algorithmicdo}{\textbf{hacer}}
\renewcommand{\algorithmicendfor}{\algorithmicend\ \algorithmicfor}
\renewcommand{\algorithmicwhile}{\textbf{mientras}}
\renewcommand{\algorithmicendwhile}{\algorithmicend\ \algorithmicwhile}
\renewcommand{\algorithmicloop}{\textbf{repetir}}
\renewcommand{\algorithmicendloop}{\algorithmicend\ \algorithmicloop}
\renewcommand{\algorithmicrepeat}{\textbf{repetir}}
\renewcommand{\algorithmicuntil}{\textbf{hasta que}}
\renewcommand{\algorithmicprint}{\textbf{imprimir}} 
\renewcommand{\algorithmicreturn}{\textbf{devolver}} 
\renewcommand{\algorithmictrue}{\textbf{cierto }} 
\renewcommand{\algorithmicfalse}{\textbf{falso }} 
 % mi archivo de traducción
\usepackage{titling} %Para llamar los datos
\usepackage{titlesec} 
\usepackage{lipsum} %Muestras de Lorem Ipsum
\usepackage{tikz}
\usepackage{afterpage}
\newcommand\blankpage{%
	\null
	\thispagestyle{empty}%
	\addtocounter{page}{0}%
	\newpage}

\usepackage{acronym} % Manejo de acronimos

\usepackage{multicol} % Uso de split de columnas para las tablas
\usepackage{multirow} % Uso de split de filas para las tablas
\usepackage{colortbl} % Coloracion de celdas
\usepackage{graphicx}
\usepackage{caption}

%%%%%%%%%%%%%%%%%%%%%%%%%%%%%%%%%%%%%%%%%%%%%%%%%%
%%% Seccion de escritura de codigo
\usepackage{listings}
\renewcommand{\lstlistingname}{Código}
\renewcommand\lstlistlistingname{Índice de Código}
\lstdefinestyle{custombash}{
	belowcaptionskip=1\baselineskip,
	breaklines=true,
	frame=L,
	xleftmargin=\parindent,
	language=bash,
	showstringspaces=false,
	basicstyle=\footnotesize\ttfamily,
	keywordstyle=\bfseries\color{green!40!black},
	commentstyle=\itshape\color{purple!40!black},
	identifierstyle=\color{black},
	stringstyle=\color{orange},
	captionpos=b
}
\lstset{escapechar=@,style=custombash}

%%%%%%%%%%%%%%%%%%%%%%%%%%%%%%%%%%%%%%%%%%%%%%%%%%
%%%%%%%%%%%%%%%%%%
%%%% Datos de la universidad, tesista y asesor
\title{{Esquema de Monitorización en Seguridad para entorno IoT }}
\author{Briggette Olenka Román Huaytalla}
\def\supervisor{Mg. Gipsy Miguel Arrunategui Angulo}
\def\university{\href{http://www.uni.edu.pe/}{Universidad Nacional de Ingeniería}}
\def\faculty{Facultad de Ciencias}
\def\school{Escuela Profesional de Ciencia de la Computación}

%%%%%%%%%%%%%%%%%%%
%%%%%%%%%%%%%%%%%%%
%%%%Inicio de archivo
\begin{document}

%%%%%%%%%%%%%%%%%%%
%%% Caratula

	\begin{titlepage}
		\begin{center}
		
		{\scshape \huge \university \par}
		\vspace{0.2cm}
		{\scshape \Large \faculty \par}
		\vspace{0.2cm}
		{\scshape \large \school \par}
%		\vspace{0.2cm}

		\begin{figure}[h]
			\centering
			\includegraphics[scale = 0.5]{figure/log_uni}
		\end{figure}
%		\vspace{0.2cm}
		
		{\LARGE \thetitle}\\[1cm]
		{\Large \textbf{SEMINARIO DE TESIS II}}\\[0.5cm]
		\large\emph{Autor: \\[0.3cm]}
		{\theauthor}\\[0.5cm]
		\large\emph{Asesor:\\[0.3cm] }
		{\supervisor }\\[1cm]
		{\large Diciembre, 2018}\\[4cm] 
		\end{center}
			
	\end{titlepage}
\afterpage{\blankpage}
%%%%%%%%%%%%%%%%%%%%

\frontmatter

%%%%%%%%%%%%%%%%%%%%

%%%%%%%%%%%%%%%%%%%%%%%
%%% Estilo de capitulos

\titleformat{\chapter}[display]
{\Large\itshape}
{\filright\MakeUppercase{\chaptertitlename} }
{1ex}
{\titlerule\vspace{1ex}\filleft}
[\vspace{1ex}\titlerule]

%%%%%%%%%%%%%%%%%%%%%%%

%%%%%%%%%%%%%%%%%%%%%%%
%%% Resumen

\chapter{Resumen}
	
		\lipsum[1] 
	
	\vspace{10px}	
	\textbf{PALABRAS CLAVE:} IoT, Ciberseguridad, OWASP, vulnerabilidades, Gestión de Riesgos.
	
	\vspace{100px}
	
	\textit{El hardware es fácil de proteger: bloquearlo en una habitación, asegurarlo en un escritorio, o comprar uno de repuesto. La información posee más de un problema. Puede existir en más de un lugar, ser transportada a medio planeta en segundos, y ser robada sin que uno tenga conocimiento.}
	
	\begin{flushright}
		(Bruce Schneider, Protege tu Macintosh, 1994)
	\end{flushright}

%%%%%%%%%%%%%%%%%%%%%%%%%
%%% Indices y Listas

\afterpage{\blankpage}
\tableofcontents
\afterpage{\blankpage}
\listoffigures
\afterpage{\blankpage}
\listoftables
\afterpage{\blankpage}
\lstlistoflistings
\afterpage{\blankpage}

%%%%%%%%%%%%%%%%%%%%%%%%%%

%%%%%%%%%%%%%%%%%%%%%%%%%%
%%%Lista de acronimos

\newpage
\chapter*{Lista de Acrónimos}
\begin{acronym}[MPC] % Give the longest label here so that the list is nicely aligned

\acro{OWASP}{Open Web Application Security Project} \\
\acro{IoT}{Internet of Things}\\
\acro{RPi}{Raspberry Pi} \\
\end{acronym}
\afterpage{\blankpage}

%%%%%%%%%%%%%%%%%%%%%%%%%%

%%%%%%%%%%%%%%%%%%%%%%%%%%
%%% Agradecimientos

\chapter*{Agradecimientos}
\lipsum[1]
\afterpage{\blankpage}

%%%%%%%%%%%%%%%%%%%%%%%%%%
%%%%%%%%%%%%%%%%%%%%%%%%%%
%%% Formato de capitulos
	
\titleformat{\chapter}[display]
{\bfseries\Large}
{\filright\MakeUppercase{\chaptertitlename} \Huge\thechapter}
{1ex}
{\titlerule\vspace{1ex}\filleft}
[\vspace{1ex}\titlerule]	
%%%%%%%%%%%%%%%%%%%%%%%%%%

%%%%%%%%%%%%%%%%%%%%%%%%%%
\mainmatter
%%%%%%%%%%%%%%%%%%%%%%%%%%

%%%%%%%%%%%%%%%%%%%%%%%%%%
%%% Capitulos
\chapter{Introducción}
\lipsum
\section{title}
\subsection{title}
\lipsum	
\afterpage{\blankpage}
\chapter{Estado del arte}
\lipsum[1-2]
adasdad \ac{OWASP}

\section{title}
\lipsum[1-3]
\subsection{title}
\lipsum[1] \cite{ISO27005}\cite{NISTGaithersburg2012}.

\afterpage{\blankpage}
\chapter{Implementación}

\begin{lstlisting}[language=bash, caption= Configuración de red del RPi]
# Listamos las redes para encontrar a la que deseamos conectarnos
$ sudo iwlist wlan0 scan

#Agregamos los datos de la red a conectarse
$ sudo nano /etc/wpa_supplicant/wpa_supplicant.conf



# Agregar lo siguiente dentro del archivo wpa_supplicant.conf
network={
ssid="wifiName"
psk="wifiPassword"
}

$ sudo reboot	
\end{lstlisting}
\afterpage{\blankpage}
\chapter{Resultados y discusiones}
\lipsum[1]
\begin{table}[h!]
	\begin{center}
		\caption{Uso de split y color de celdas.}
		\label{tab:table1}
		\begin{tabular}{l|p{3cm}|r}
			
			\textbf{Value 1} & \textbf{Value 2} & \textbf{Value 3}\\
			$\alpha$ & $\beta$ & $\gamma$ \\
			\hline
			
			\multicolumn{2}{c|}{\cellcolor[gray]{.9}} & a\\
			\multicolumn{2}{c|}{\multirow{-2}{*}{\cellcolor[gray]{.9}1234}} & b\\
			\hline
			3 & \centering 23.113231 & c\\
			4 & 25.113231 & d\\
		\end{tabular}
	\end{center}
\end{table}
\afterpage{\blankpage}
\chapter{Conclusiones}
\section{Trabajo a futuro}
\section{Conclusiones}

%%%%%%%%%%%%%%%%%%%%%%%%%%

%%%%%%%%%%%%%%%%%%%%%%%%%%
%%% Bibliografia
\bibliography{reference}
\bibliographystyle{unsrt}
\afterpage{\blankpage}
%%%%%%%%%%%%%%%%%%%%%%%%%%

%%%%%%%%%%%%%%%%%%%%%%%%%%
%%% Apendices
\appendix
\chapter{Nombre del apéndice} % Main appendix title
\label{AppendA}
%%%%%%%%%%%%%%%%%%%%%%%%%%
	
\end{document}